\section*{第 \ref{ch:b1:geological-maps} 章——地质图}

\begin{solution}{exo:b1:geological-maps:1}
\omterm{def:b1:geological-maps:attitude}{走向}：平面水平线的\omterm{def:b1:spherical-astronomy:horizontal}{方位角}；\omterm{def:b1:geological-maps:attitude}{倾角}：平面与水平面之间的夹角，
沿垂直于\omterm{def:b1:geological-maps:attitude}{走向}的方向量取；\omterm{def:b1:geological-maps:attitude}{倾向}：
沿平面最陡下倾线的\omterm{def:b1:spherical-astronomy:horizontal}{方位角}。对于 $135/40$，按右手定则，\omterm{def:b1:geological-maps:attitude}{走向}为
$135 - 90 = \ang{045}$。符号：\omterm{def:b1:geological-maps:attitude}{走向}线朝向
N\,45\,E，在其东南侧（朝向 \ang{135}）画一短刻线，标注 40。
\end{solution}

\begin{solution}{exo:b1:geological-maps:2}
$z = 500 - \tan\ang{30}\,(x\sin\ang{90} + y\cos\ang{90}) = 500 - 0.577\,x$。
在 $(1000, 0)$ 处：$500 - 577 = \qty{-77}{m}$。在 $(0, 1000)$ 处：\qty{500}{m}，
该平面沿南北向\omterm{def:b1:geological-maps:attitude}{走向}呈水平。
\end{solution}

\begin{solution}{exo:b1:geological-maps:3}
沿等高线延伸的迹线为水平接触面，高程为 \qty{350}{m}。
横穿丘陵和山谷、呈直线延伸的迹线属于
直立接触面。
\end{solution}

\begin{solution}{exo:b1:geological-maps:4}
沿河流，\omterm{def:b1:geological-maps:fold}{背斜}核部出露片岩和花岗岩
（位于砂岩之下）；\omterm{def:b1:geological-maps:fold}{向斜}核部为白垩。
最老的岩层位于序列的最下部：岩层向上拱起
且侵蚀切割最深处，最老的地层出露；岩层下沉处，
最新的地层保存在中部。
\end{solution}

\begin{solution}{exo:b1:geological-maps:5}
前两点高程相同：直线 $y = 0$ 为\omterm{def:b1:geological-maps:attitude}{走向}线，
因此\omterm{def:b1:geological-maps:attitude}{走向}为东西向。接触面低于 \qty{100}{m}，
位于北面 \qty{500}{m} 处：它向北倾，$\tan\delta = 100/500 = 0.2$，
$\delta \approx \ang{11.3}$；\omterm{def:b1:geological-maps:attitude}{倾向} \ang{000}，\omterm{def:b1:geological-maps:attitude}{走向} \ang{270}（按右手定则）。
用定理验算：$\overrightarrow{P_1P_2} = (1000, 0, 0)$，
$\overrightarrow{P_1P_3} = (0, 500, -100)$，$\vec n = (0, 100\,000,
500\,000)$，$\cos\delta = 500\,000/509\,902 \approx 0.981$，$\delta \approx
\ang{11.3}$，且 $(n_x, n_y) = (0, 100\,000)$ 指向北，结果相同。
\end{solution}

\begin{solution}{exo:b1:geological-maps:6}
下游方向为北，$s = 0.03$（\ang{1.72}）。(a) $a = 0 < s$：V 形指向上游
（南），沿等高线。(b) $a = -\tan\ang{20} < 0 < s$：V 形指向上游。
(c) $a = \tan\ang{1} \approx 0.017 < s$：V 形指向上游。(d) $a = \tan\ang{20}
\approx 0.364 > s$：V 形指向下游（北）。
\end{solution}

\begin{solution}{exo:b1:geological-maps:7}
$\xi_{\text{谷底}} - \xi_{\text{侧面}} = h/(a - s) = 80/(0.364 - 0.03)
\approx \qty{240}{m}$：侧面迹线位于顶点上游约
\qty{240}{m} 处。
\end{solution}

\begin{solution}{exo:b1:geological-maps:8}
\omterm{def:b1:geological-maps:fold}{背斜}的闭合点指向倾伏方向：轴线
向北倾伏。轴线北倾时，\omterm{def:b1:geological-maps:fold}{向斜}的\omterm{def:b1:geological-maps:geological-map}{露头}在
南侧闭合。
\end{solution}

\begin{solution}{exo:b1:geological-maps:9}
东南偏东一侧出露较新的岩层（白垩和页岩，与
对侧较老的岩层相接）：该断块向下移动。断层符号的刻线
位于该侧，即下降盘一侧。
\end{solution}

\begin{solution}{exo:b1:geological-maps:10}
两平面的公共直线垂直于两个法向量，因此
平行于 $\vec n_1 \times \vec n_2$。对于 $090/30$：$\vec n_1 = (0.577, 0,
1)$；对于 $270/30$：$\vec n_2 = (-0.577, 0, 1)$。于是 $\vec n_1 \times \vec n_2
= (0\cdot 1 - 1\cdot 0,\ 1\cdot(-0.577) - 0.577\cdot 1,\ 0) = (0, -1.155,
0)$：一条水平的南北向直线。第一平面向东倾，第二
平面向西倾：两者背离该直线倾斜，这是一个
水平轴\omterm{def:b1:geological-maps:fold}{背斜}的脊线。
\end{solution}

\begin{solution}{exo:b1:geological-maps:11}
直立岩墙的迹线是过原点、\omterm{def:b1:geological-maps:attitude}{走向}为
\ang{045} 的直线，与高程无关：在山丘上它也位于
同一直线上。倾斜岩墙为平面 $z = -\tan\ang{60}\,d$，其中 $d$ 为
到过原点\omterm{def:b1:geological-maps:attitude}{走向}线的水平距离，
朝东南为正。在 $z = \qty{300}{m}$ 处：$d = -300/\tan\ang{60}
\approx \qty{-173}{m}$：迹线位于过原点\omterm{def:b1:geological-maps:attitude}{走向}线的西北约
\qty{173}{m} 处，位于上倾一侧。
\end{solution}

\begin{solution}{exo:b1:geological-maps:12}
\[
  \frac{\partial z}{\partial x} \approx \frac{405 - 420}{100} = -0.15, \qquad
  \frac{\partial z}{\partial y} \approx \frac{412 - 420}{100} = -0.08,
\]
因此坡度为 $\arctan\sqrt{0.15^2 + 0.08^2} = \arctan 0.17 \approx
\ang{9.6}$。地面
沿 $-\nabla z = (0.15, 0.08)$ 下降，其\omterm{def:b1:spherical-astronomy:horizontal}{方位角}为
$\arctan(0.15/0.08) \approx \ang{62}$：该环流圈朝向东北偏东。
\end{solution}

\begin{solution}{pb:b1:geological-maps:1}
\textbf{1.} 片岩、花岗岩、砂岩、石灰岩、页岩、白垩、冲积层。

\textbf{2.} 由西向东：白垩和页岩、石灰岩、砂岩，然后
在 $x = \qty{3.5}{km}$ 附近为片岩和花岗岩，之后又为砂岩、石灰岩和页岩，
最后为白垩；谷底全程为冲积层。

\textbf{3.} 砂岩，其下沿河流为片岩和花岗岩：
核部为最老的岩层，属\omterm{def:b1:geological-maps:fold}{背斜}。

\textbf{4.} 白垩，为最新的岩层，属\omterm{def:b1:geological-maps:fold}{向斜}。

\textbf{5.} 砂岩底面在\omterm{def:b1:geological-maps:fold}{背斜}轴部最高，
只有沿河流处侵蚀才切割至其下方。

\textbf{6.} 不整合：砂岩覆于两种不同的较老岩层之上，
即片岩与花岗岩，并截断了花岗岩的轮廓。

\textbf{7.} $435 - 261 = 174$, $432 - 428 = 4$, $512 - 586 = \qty{-74}{m}$.

\textbf{8.} $\overrightarrow{E_1E_2} = (600, 1200, -170)$,
$\overrightarrow{E_1E_3} = (1100, 400, -248)$.

\textbf{9.} 叉积为
\[
  (-229\,600,\ -38\,200,\ -1\,080\,000);
\]
取向朝上，$\vec n_E = (229\,600,\ 38\,200,\ 1\,080\,000)$。

\textbf{10.} $\|\vec n_E\| \approx 1\,104\,800$; $\cos\delta \approx 0.9776$,
$\delta \approx \ang{12.2}$.

\textbf{11.} $(229\,600, 38\,200)$ 的\omterm{def:b1:spherical-astronomy:horizontal}{方位角}：$\arctan(229\,600/38\,200)
\approx \ang{80.6}$；\omterm{def:b1:geological-maps:attitude}{产状} $081/12$，\omterm{def:b1:geological-maps:attitude}{走向} $\approx \ang{351}$。

\textbf{12.} 是：东翼上的符号，其\omterm{def:b1:geological-maps:attitude}{走向}线接近
南北向，东侧带有\omterm{def:b1:geological-maps:attitude}{倾向}短线，标注为 12。

\textbf{13.} $524 - 669 = -145$, $421 - 459 = -38$, $319 - 209 =
\qty{110}{m}$.

\textbf{14.} $\overrightarrow{W_1W_2} = (700, 1200, 107)$ 和
$\overrightarrow{W_1W_3} = (1100, -600, 255)$。二者的叉积，
取向朝上，为
\[
  (370\,200,\ -60\,800,\ -1\,740\,000) \quad\longrightarrow\quad
  \vec n_W = (-370\,200,\ 60\,800,\ 1\,740\,000).
\]

\textbf{15.} $\cos\delta = 1\,740\,000/\|\vec n_W\| \approx 0.9776$, $\delta
\approx \ang{12.2}$；\omterm{def:b1:geological-maps:attitude}{倾向}为 $(-370\,200, 60\,800)$ 的\omterm{def:b1:spherical-astronomy:horizontal}{方位角}，
$\approx \ang{279.3}$。\omterm{def:b1:geological-maps:attitude}{产状} $279/12$，\omterm{def:b1:geological-maps:attitude}{走向} $\approx \ang{189}$：
西翼向西北西倾，东翼向东北东倾。

\textbf{16.} $\vec n_E \times \vec n_W = (804 \times 10^6,\ -799\,320 \times
10^6,\ 28\,101 \times 10^6)$（取整）。

\textbf{17.} 向下方向为 $(-804, 799\,320, -28\,101) \times
10^6$：倾伏向 $\arctan(-804/799\,320) \approx \ang{0}$，即正北；倾伏角
$\arctan\bigl(28\,101/\sqrt{804^2 + 799\,320^2}\bigr) \approx \ang{2.0}$。\omterm{def:b1:geological-maps:fold}{褶皱
轴}向北倾伏 \ang{2}。

\textbf{18.} $8000 \times \tan\ang{2.0} \approx \qty{280}{m}$ 低，
约在 \qty{-50}{m} 处。\omterm{def:b1:geological-maps:fold}{背斜}核部在河流附近最高，向北下沉，
因此其\omterm{def:b1:geological-maps:geological-map}{露头}向北闭合，与图所示一致。

\textbf{19.} $\overrightarrow{F_1F_2} = (1627, 2500, -276)$ 和
$\overrightarrow{F_1F_3} = (3124, 5500, 77)$ 给出朝上的法向量
\[
  (1\,710\,500,\ -987\,503,\ 1\,138\,500), \qquad
  \delta = \arccos\frac{1\,138\,500}{2\,279\,700} \approx \ang{60},
\]
并得\omterm{def:b1:geological-maps:attitude}{倾向} $\approx \ang{120}$。\omterm{def:b1:geological-maps:attitude}{产状} $120/60$，\omterm{def:b1:geological-maps:attitude}{走向}
\ang{030}。

\textbf{20.} 断层向东南东倾，故该侧为
上盘，即位于倾斜断层面之上的块体；该块体向下移动（地表为较新的岩层）：正断层。

\textbf{21.} 断层面向东南东倾，顺河流方向看，其倾斜远比河床下降更陡：
较低的地面与断层面相遇于更远的下倾方向，
因此断层迹线在山谷中向东南东偏折——即
对断层面应用 V 形法则。

\textbf{22.} 花岗岩侵入；砂岩沉积前发生侵蚀；
白垩沉积；\omterm{def:b1:geological-maps:fold}{褶皱}；断层作用；冲积层沉积。

\textbf{23.} 晚于 9000 万年前（白垩已\omterm{def:b1:geological-maps:fold}{褶皱}），且早于
258 万年前（冲积层未被断层错断）；\omterm{def:b1:geological-maps:fold}{褶皱}
先发生，断层作用随后。

\textbf{24.} 断层自 F1（\qty{405}{m}）下降 \qty{1405}{m}，至
\qty{-1000}{m}；\omterm{def:b1:geological-maps:attitude}{倾角}为 \ang{60} 时，对应的水平距离为
$1405/\tan\ang{60} \approx \qty{811}{m}$（沿\omterm{def:b1:geological-maps:attitude}{倾向}）。水井应钻于
约 \textbf{\qty[text-series-to-math=true]{810}{m}，位于 F1 东南东方向}（\omterm{def:b1:spherical-astronomy:horizontal}{方位角} \ang{120}）。
\end{solution}
